% ============================================================================
% LANGENTWURF: Repaso Unidad 5 - La ropa y los adjetivos
% ============================================================================
% Klasse 8 | 45 Minuten (Einzelstunde) | Tier 1 (vollständig)
% Fokus: Kleidungsvokabeln + Adjektive zur Beschreibung
% ============================================================================

\documentclass[11pt,a4paper]{article}

% ============================================================================
% PACKAGES
% ============================================================================
\usepackage[utf8]{inputenc}
\usepackage[T1]{fontenc}
\usepackage[ngerman]{babel}
\usepackage{geometry}
\usepackage{fancyhdr}
\usepackage{lastpage}
\usepackage{xcolor}
\usepackage{tcolorbox}
\usepackage{enumitem}
\usepackage{tabularx}
\usepackage{longtable}
\usepackage{array}
\usepackage{booktabs}
\usepackage{hyperref}
\usepackage{amssymb}

% ============================================================================
% KONFIGURATION
% ============================================================================

\newcommand{\plantier}{1}
\newcommand{\fachid}{spanisch}
\newcommand{\fach}{Spanisch}
\newcommand{\fachfarbe}{c41e3a}

% --- Metadaten ---
\newcommand{\jahrgangsstufe}{8}
\newcommand{\schulart}{Gesamtschule}
\newcommand{\stundenthema}{Repaso: La ropa y los adjetivos}
\newcommand{\reihenthema}{Unidad 5 -- Mi ropa preferida}
\newcommand{\stundennummer}{Repaso}
\newcommand{\unterrichtsdatum}{\underline{\hspace{3cm}}}
\newcommand{\unterrichtszeit}{45 Minuten (Einzelstunde)}

% --- Lerngruppe ---
\newcommand{\lerngruppengroesse}{ca. 24}
\newcommand{\gerniveau}{A1}
\newcommand{\lehrwerk}{Qué pasa 1}
\newcommand{\lehrwerkseite}{S. 88--101, bes. S. 101 (Gramática y comunicación)}

% ============================================================================
% LAYOUT
% ============================================================================
\geometry{left=2.5cm, right=2.5cm, top=2.5cm, bottom=2.5cm}
\definecolor{fach}{HTML}{\fachfarbe}
\definecolor{gray}{HTML}{666666}
\definecolor{lightgray}{HTML}{f5f5f5}
\definecolor{successgreen}{HTML}{2a7a4b}
\definecolor{warningorange}{HTML}{b35c00}
\definecolor{basisgreen}{HTML}{2E7D32}
\definecolor{standardblue}{HTML}{1565C0}
\definecolor{erweiterungpurple}{HTML}{7B1FA2}

\tcbuselibrary{skins,breakable}

\newtcolorbox{infobox}[1][]{
    colback=lightgray,
    colframe=fach,
    fonttitle=\bfseries,
    title=#1,
    breakable
}

\newtcolorbox{scorebox}{
    colback=white,
    colframe=fach,
    fonttitle=\bfseries,
    title=Didaktik-Score,
    breakable
}

\newtcolorbox{tippbox}{
    colback=white,
    colframe=warningorange,
    boxrule=1pt,
    left=5pt, right=5pt, top=3pt, bottom=3pt
}

% Differenzierungsmarker
\newcommand{\basis}{\textcolor{basisgreen}{\textbf{[B]}}}
\newcommand{\standard}{\textcolor{standardblue}{\textbf{[S]}}}
\newcommand{\erweiterung}{\textcolor{erweiterungpurple}{\textbf{[E]}}}

% Kopf-/Fußzeile
\pagestyle{fancy}
\fancyhf{}
\fancyhead[L]{\small\textcolor{gray}{\fach{} -- Kl. \jahrgangsstufe{} -- \stundenthema}}
\fancyhead[R]{\small\textcolor{gray}{Langentwurf Tier 1}}
\fancyfoot[C]{\small\thepage{} / \pageref{LastPage}}
\renewcommand{\headrulewidth}{0.4pt}
\renewcommand{\footrulewidth}{0pt}

% ============================================================================
% DOKUMENT
% ============================================================================
\begin{document}

% ============================================================================
% DECKBLATT
% ============================================================================
\begin{titlepage}
    \centering
    \vspace*{2cm}

    {\large\textcolor{gray}{\schulart{} -- Klasse \jahrgangsstufe}}\\[2cm]

    {\Huge\textcolor{fach}{\textbf{Langentwurf}}}\\[0.5cm]
    {\large Tier 1 -- Vollständig}\\[2cm]

    {\LARGE\textbf{\stundenthema}}\\[0.5cm]
    {\large im Rahmen der Unterrichtsreihe}\\[0.3cm]
    {\Large\textit{\reihenthema}}\\[2cm]

    \begin{tabular}{rl}
        \textbf{Fach:} & \fach{} (\gerniveau) \\[0.3cm]
        \textbf{Klasse:} & \jahrgangsstufe \\[0.3cm]
        \textbf{Lehrwerk:} & \lehrwerk, \lehrwerkseite \\[0.3cm]
        \textbf{Dauer:} & \textbf{45 Minuten (Einzelstunde)} \\[0.3cm]
        \textbf{Datum:} & \unterrichtsdatum \\[0.3cm]
        \textbf{Format:} & Papierbasiert (keine Tablets) \\[0.3cm]
    \end{tabular}

    \vfill

    {\normalsize\textcolor{gray}{Repaso-Stunde: Wiederholung und Festigung}}\\[1cm]

\end{titlepage}

% ============================================================================
% INHALTSVERZEICHNIS
% ============================================================================
\tableofcontents
\newpage

% ============================================================================
% 1. BEDINGUNGSANALYSE
% ============================================================================
\section{Bedingungsanalyse}

\subsection{Lerngruppe}
\begin{tabular}{ll}
    Kursgröße: & \lerngruppengroesse{} SuS \\
    Jahrgangsstufe: & \jahrgangsstufe \\
    Sprachniveau: & \gerniveau{} (GER) \\
    Fremdsprache: & 2. Fremdsprache (seit Klasse 7) \\
\end{tabular}

\subsection{Voraussetzungen}
Die SuS haben in den vorangegangenen Stunden der Unidad 5 bereits behandelt:
\begin{itemize}
    \item Kleidungsvokabeln (Grundwortschatz)
    \item Farbadjektive und deren Genus-Anpassung
    \item Modalverben \textit{querer, poder, tener que}
    \item Zahlen 100--1000
    \item Einkaufsdialoge (Grundstruktur)
\end{itemize}

\textbf{Letzte Stunde (Repaso Números):} Fokus auf Zahlen 100--1000. Die vorliegende Stunde ergänzt dies durch den Fokus auf \textbf{Kleidungsbeschreibungen mit Adjektiven}.

\subsection{Differenzierung (intern)}

\textbf{WICHTIG:} Keine Sternchen oder Niveaumarkierungen auf Schülermaterial!

\begin{tabular}{|l|p{3.5cm}|p{3.5cm}|p{3.5cm}|}
\hline
& \textbf{\basis{} Basis} & \textbf{\standard{} Standard} & \textbf{\erweiterung{} Erweiterung} \\
\hline
\textbf{Ziel} & Berufsreife & Mittlere Reife & Gymnasium \\
\hline
\textbf{Inhalt} & Zuordnen, Lücken mit Hilfe & Eigenständiges Einsetzen & Freie Beschreibung \\
\hline
\textbf{Hilfen} & Wortliste, Beispielsätze & Teils Beispiele & Nur Impulse \\
\hline
\end{tabular}

\subsection{Besonderheiten der Stunde}
\begin{itemize}
    \item \textbf{Nur 45 Minuten} (Einzelstunde statt Doppelstunde)
    \item \textbf{Papierbasiert} -- keine digitalen Geräte
    \item \textbf{Repaso-Charakter} -- Wiederholung, keine Neueinführung
    \item \textbf{Mündliche Übung begrenzt} -- kurze Partnerphase nach Vokabelrevision
\end{itemize}

% ============================================================================
% 2. SACHANALYSE
% ============================================================================
\section{Sachanalyse}

\subsection{Sprachliche Strukturen}

\subsubsection{Kleidungsvokabular}
\begin{tabular}{ll@{\hspace{2cm}}ll}
\textbf{Spanisch} & \textbf{Deutsch} & \textbf{Spanisch} & \textbf{Deutsch} \\
\midrule
la camiseta & das T-Shirt & la camisa & das Hemd \\
el jersey & der Pullover & la sudadera & der Kapuzenpulli \\
los pantalones & die Hose & el vestido & das Kleid \\
la falda & der Rock & la chaqueta & die Jacke \\
el abrigo & der Mantel & los zapatos & die Schuhe \\
las botas & die Stiefel & las zapatillas & die Turnschuhe \\
\end{tabular}

\subsubsection{Beschreibende Adjektive (Fokus dieser Stunde)}
\begin{tabular}{ll@{\hspace{2cm}}ll}
\textbf{Spanisch} & \textbf{Deutsch} & \textbf{Spanisch} & \textbf{Deutsch} \\
\midrule
bonito/a & hübsch & moderno/a & modern \\
cómodo/a & bequem & elegante & elegant \\
ancho/a & weit & estrecho/a & eng \\
largo/a & lang & corto/a & kurz \\
deportivo/a & sportlich & sencillo/a & schlicht \\
\end{tabular}

\subsubsection{Grammatische Struktur: Adjektiv-Konkordanz}
\begin{itemize}
    \item Adjektive auf \textbf{-o/-a} passen sich dem Genus an:\\
    \textit{el jersey \textbf{bonito}} / \textit{la camiseta \textbf{bonita}}
    \item Adjektive auf \textbf{-e} oder Konsonant bleiben unverändert:\\
    \textit{el vestido \textbf{elegante}} / \textit{la falda \textbf{elegante}}
    \item Plural: \textit{los pantalones \textbf{cómodos}} / \textit{las botas \textbf{cómodas}}
\end{itemize}

\subsection{Redemittel (S. 101 ``Über Kleidung sprechen'')}
\begin{itemize}
    \item \textit{El/La [Kleidung] es muy [Adjektiv].}
    \item \textit{Me gusta el/la [Kleidung] porque es [Adjektiv].}
    \item \textit{El/La [Kleidung] es demasiado [grande/pequeño/caro].}
    \item \textit{Llevo un/una [Kleidung] [Adjektiv].}
\end{itemize}

\subsection{Kontrastive Betrachtung}
\begin{itemize}
    \item \textbf{Positiver Transfer:} Viele Adjektive sind transparent (moderno, elegante)
    \item \textbf{Interferenz:} Im Deutschen stehen Adjektive vor dem Nomen, im Spanischen meist danach
    \item \textbf{Stolperstein:} Genus-Anpassung bei -o/-a vergessen (typischer Fehler)
\end{itemize}

% ============================================================================
% 3. DIDAKTISCHE ANALYSE
% ============================================================================
\section{Didaktische Analyse}

\subsection{Stellung in der Sequenz}
Dies ist eine \textbf{Repaso-Stunde} zum Abschluss der Unidad 5 ``Mi ropa preferida''. Nach der vorherigen Repaso-Stunde (Fokus: Zahlen 100--1000) liegt hier der Schwerpunkt auf \textbf{Kleidungsbeschreibungen mit Adjektiven}.

\subsection{Kompetenzziele (Rahmenplan MV)}
\begin{itemize}
    \item \textbf{Kommunikativ:} Sprechen (monologisch/dialogisch) -- Kleidung beschreiben
    \item \textbf{Sprachlich:} Verfügen über sprachliche Mittel (Wortschatz, Grammatik)
    \item \textbf{Methodisch:} Wortschatzarbeit, Strukturieren
\end{itemize}

\subsection{Legitimation}
\textbf{Rahmenplan Spanisch MV (Kl. 7--10):}
\begin{itemize}
    \item Kompetenzbereich: Sprechen -- zusammenhängendes Sprechen
    \item Standard: ``Gegenstände und Personen beschreiben''
    \item Themenfeld: ``Alltagsleben'' (Kleidung, Einkaufen)
\end{itemize}

\subsection{Didaktische Reduktion}
\begin{itemize}
    \item \textbf{Fokus auf 10 Adjektive} (nicht alle aus Unidad 5)
    \item \textbf{Nur Singular-Formen} im Hauptteil (Plural nur bei Extension)
    \item \textbf{Keine Steigerung} (más bonito) -- erst später
\end{itemize}

% ============================================================================
% 4. STUNDENZIELE
% ============================================================================
\section{Stundenziele}

\subsection{Grobziel}
Die SuS erweitern ihre \textbf{kommunikative Kompetenz} im Bereich Sprechen und Schreiben, indem sie Kleidungsstücke mit passenden Adjektiven beschreiben und dabei die Genus-Konkordanz korrekt anwenden.

\subsection{Feinziele (operationalisiert)}
\begin{tabular}{|c|p{7cm}|c|c|}
\hline
\textbf{Nr.} & \textbf{Die SuS können...} & \textbf{AFB} & \textbf{Diff.} \\
\hline
FZ1 & die 12 Kleidungsvokabeln korrekt zuordnen und aussprechen & I & alle \\
\hline
FZ2 & die 10 beschreibenden Adjektive mit korrekter Bedeutung verwenden & I & alle \\
\hline
FZ3 & Adjektive mit korrekter Genus-Endung (-o/-a/-e) einsetzen & II & alle \\
\hline
FZ4 & vollständige Beschreibungssätze nach Muster bilden & II & \standard{}/\erweiterung{} \\
\hline
FZ5 & einem Partner mündlich 2--3 Kleidungsstücke beschreiben & II & \standard{}/\erweiterung{} \\
\hline
FZ6 & eigene kreative Beschreibungen ohne Vorgabe formulieren & III & \erweiterung{} \\
\hline
\end{tabular}

% ============================================================================
% 5. METHODISCHE ANALYSE
% ============================================================================
\section{Methodische Analyse}

\subsection{Phasierung (45 Minuten Einzelstunde)}
\begin{enumerate}
    \item \textbf{Einstieg (5 Min):} Bildimpuls -- Kleidung benennen
    \item \textbf{Vokabelrevision (8 Min):} LB durchgehen, Adjektive einführen/wiederholen
    \item \textbf{Übung 1 (10 Min):} AB -- Adjektive einsetzen (Lückentext)
    \item \textbf{Übung 2 (10 Min):} AB -- Beschreibungssätze bilden
    \item \textbf{Partnerübung (7 Min):} Mündliche Beschreibung (nach Vokabelrevision!)
    \item \textbf{Sicherung (5 Min):} Zusammenfassung, Lösungen
\end{enumerate}

\subsection{Sozialformen}
\begin{itemize}
    \item \textbf{Plenum:} Einstieg, Vokabelrevision, Sicherung
    \item \textbf{Einzelarbeit:} AB-Übungen
    \item \textbf{Partnerarbeit:} Kurze mündliche Übung (bewusst begrenzt)
\end{itemize}

\subsection{Medien}
\begin{itemize}
    \item \textbf{LB-05-repaso-ropa.pdf} -- Lernblatt (``Textbuch-Ersatz'')
    \item \textbf{AB-05-repaso-ropa.pdf} -- Arbeitsblatt mit Übungen
    \item Tafel/Whiteboard für Beispielsätze
    \item Optional: Bildkarten Kleidung
\end{itemize}

% ============================================================================
% 6. VERLAUFSPLANUNG
% ============================================================================
\section{Verlaufsplanung}

\textbf{Einzelstunde (45 Minuten)}

\begin{longtable}{|p{1.2cm}|p{1.8cm}|p{4.5cm}|p{3.5cm}|p{1.2cm}|p{1.5cm}|}
\hline
\textbf{Zeit} & \textbf{Phase} & \textbf{Lehrerhandeln} & \textbf{Schülerhandeln} & \textbf{SF} & \textbf{Medien} \\
\hline
\endhead

0--5 & Einstieg &
Begrüßung auf Spanisch. Zeigt Bild/Zeichnung einer Person mit verschiedenen Kleidungsstücken. Fragt: \textit{``¿Qué lleva?''} &
Benennen Kleidungsstücke spontan (aktivieren Vorwissen) &
Plenum &
Tafel/ Bild \\
\hline

5--13 & Vokabel-revision &
Teilt LB aus. Geht Kleidungsvokabeln durch (chorisches Sprechen). Führt dann die 10 Adjektive ein: \textit{``¿Cómo es el jersey? -- Es bonito, moderno...''} Erklärt Genus-Regel kurz. &
Sprechen nach, ordnen Adjektive zu, notieren ggf. auf LB &
Plenum &
LB \\
\hline

13--23 & Übung 1: Adjektive einsetzen &
Teilt AB aus. Erklärt Aufgabe 1 (Lückentext). Geht herum, unterstützt bei Bedarf. \basis{}: Hinweis auf Wortliste &
Bearbeiten AB Aufgabe 1 selbstständig. Setzen Adjektive mit korrekter Endung ein. &
EA &
AB \\
\hline

23--33 & Übung 2: Sätze bilden &
Erklärt Aufgabe 2 (Sätze nach Muster). Schreibt Beispiel an Tafel: \textit{``La camiseta es muy cómoda.''} \erweiterung{}: freie Sätze &
Bilden Beschreibungssätze. \basis{}: mit Hilfe, \standard{}: nach Muster, \erweiterung{}: frei &
EA &
AB, Tafel \\
\hline

33--40 & Partner-übung &
Leitet Partnerphase ein: \textit{``Beschreibe deinem Partner 2 Kleidungsstücke.''} Gibt Redemittel an Tafel. Achtet auf angemessene Lautstärke. &
Beschreiben sich gegenseitig Kleidungsstücke mündlich. Verwenden Redemittel. &
PA &
-- \\
\hline

40--45 & Sicherung &
Fragt nach Ergebnissen. Bespricht 2--3 Beispielsätze. Lobt korrekte Genus-Anpassung. Fasst zusammen: \textit{``Adjektive passen sich an!''} &
Nennen Beispiele, korrigieren ggf. &
Plenum &
Tafel \\
\hline

\end{longtable}

\textbf{Legende:} EA = Einzelarbeit, PA = Partnerarbeit, SF = Sozialform

% ============================================================================
% 7. ERWARTETE SCHWIERIGKEITEN
% ============================================================================
\section{Erwartete Schwierigkeiten}

\begin{tabular}{|p{4.5cm}|p{8.5cm}|}
\hline
\textbf{Schwierigkeit} & \textbf{Reaktion} \\
\hline
Genus-Endung vergessen (\textit{*la camiseta bonito}) & Kontrastives Beispiel an Tafel: \textit{el jersey bonit\textbf{o}} vs. \textit{la camiseta bonit\textbf{a}}. Merkhilfe: ``Das Adjektiv ist der Schatten des Nomens!'' \\
\hline
Verwechslung \textit{ancho/estrecho} & Gestik: Arme weit/eng. Bildliche Darstellung. \\
\hline
Adjektive auf -e: Unsicherheit ob Anpassung & Klare Regel: ``-e bleibt -e!'' Beispiel: \textit{elegante, elegante} \\
\hline
Zeitproblem (45 Min knapp) & Puffer: Partnerübung kürzen (5 statt 7 Min). Notfalls Aufgabe 2 als HA. \\
\hline
Unruhe bei Partnerarbeit & Klare Zeitvorgabe (``3 Minuten!''), Timer sichtbar, bei Bedarf abbrechen. \\
\hline
Heterogenität & Differenzierte Aufgaben auf AB: Pflicht (alle) + Extension (Schnelle) \\
\hline
\end{tabular}

% ============================================================================
% 8. DIFFERENZIERTE UNTERSTÜTZUNG
% ============================================================================
\section{Differenzierte Unterstützung}

\begin{tabular}{|l|p{4cm}|p{4cm}|p{4cm}|}
\hline
& \textbf{\basis{} Basis} & \textbf{\standard{} Standard} & \textbf{\erweiterung{} Erweiterung} \\
\hline
\textbf{Aufgabe 1} &
Adjektive aus Wortliste wählen, Beispiel vorgegeben &
Adjektive selbst einsetzen &
+ Begründung warum \\
\hline
\textbf{Aufgabe 2} &
Satzanfänge vorgegeben, nur Adjektiv ergänzen &
Nach Satzmuster arbeiten &
Freie Beschreibung ohne Muster \\
\hline
\textbf{Partnerübung} &
Redemittel ablesen erlaubt &
Redemittel als Hilfe &
Spontan sprechen \\
\hline
\textbf{Zusatz} &
-- &
Extension bei Zeit &
Kreative Aufgabe \\
\hline
\end{tabular}

% ============================================================================
% 9. TAFELANSCHRIEB
% ============================================================================
\section{Tafelanschrieb}

\begin{tcolorbox}[colback=white, colframe=black, boxrule=0.5pt]
\begin{center}
\textbf{\large La ropa y los adjetivos}
\end{center}

\vspace{0.5em}

\textbf{Adjektive -- Regeln:}
\begin{itemize}[leftmargin=*]
    \item -o/-a: passt sich an! \\
    \textit{el jersey bonit\textbf{o}} / \textit{la camiseta bonit\textbf{a}}
    \item -e: bleibt gleich! \\
    \textit{el vestido elegant\textbf{e}} / \textit{la falda elegant\textbf{e}}
\end{itemize}

\vspace{0.5em}

\textbf{Redemittel:}
\begin{itemize}[leftmargin=*]
    \item \textit{El/La ... es muy ...}
    \item \textit{Me gusta el/la ... porque es ...}
    \item \textit{Llevo un/una ... ...}
\end{itemize}
\end{tcolorbox}

% ============================================================================
% 10. HAUSAUFGABE
% ============================================================================
\section{Hausaufgabe}

\textbf{Optional} (da Repaso-Stunde):
\begin{itemize}
    \item Vokabeln wiederholen (Kleidung + Adjektive)
    \item AB fertigstellen (falls nicht geschafft)
    \item \erweiterung{}: Eigenes Outfit beschreiben (3--4 Sätze)
\end{itemize}

% ============================================================================
% 11. ANLAGEN
% ============================================================================
\section{Anlagen}

\begin{enumerate}
    \item \textbf{LB-05-repaso-ropa.pdf} -- Lernblatt (Vokabeln + Adjektive + Redemittel)
    \item \textbf{AB-05-repaso-ropa.pdf} -- Arbeitsblatt mit Übungen
    \item \textbf{ML-05-repaso-ropa.pdf} -- Musterlösung (AB mit roten Lösungen)
    \item \textbf{LH-05-repaso-ropa.pdf} -- Lehrerhinweise (Kurzplan + Lösungen)
\end{enumerate}

% ============================================================================
% 12. REFLEXIONSFRAGEN (für Lehrkraft)
% ============================================================================
\section{Reflexionsfragen}

\begin{enumerate}
    \item Haben die SuS die Genus-Anpassung verstanden und angewendet?
    \item War die Partnerübung produktiv oder zu unruhig?
    \item Reichte die Zeit für alle Aufgaben?
    \item Welche Adjektive bereiteten besondere Schwierigkeiten?
    \item War die Differenzierung ausreichend?
\end{enumerate}

% ============================================================================
% 13. DIDAKTIK-SCORE
% ============================================================================
\newpage
\section{Didaktik-Score}

\begin{scorebox}
\textbf{Bewertung der didaktischen Qualität nach 10 Dimensionen}\\
\textbf{Ziel-Score: $\geq$ 9.0 (Premium-Qualität)}

\vspace{1em}
\begin{tabular}{|c|l|c|c|p{4.5cm}|}
\hline
\textbf{\#} & \textbf{Dimension} & \textbf{Gew.} & \textbf{Score} & \textbf{Notiz} \\
\hline
1 & Differenzierung & 15\% & 9/10 & 3 Niveaus [B]/[S]/[E], qualitativ verschieden \\
\hline
2 & Taxonomie (AFB) & 10\% & 9/10 & AFB I (FZ1-2), II (FZ3-5), III (FZ6) \\
\hline
3 & Lernziel-Alignment & 10\% & 10/10 & Klar auf Rahmenplan bezogen \\
\hline
4 & Aktivierung & 10\% & 9/10 & EA + PA, konstruktives Üben \\
\hline
5 & Scaffolding & 10\% & 9/10 & Gestufte Hilfen, Redemittel \\
\hline
6 & Feedback-Qualität & 10\% & 8/10 & Mündlich im Plenum, AB-Kontrolle \\
\hline
7 & Sprachliche Klarheit & 10\% & 10/10 & Altersgerecht, klare Anweisungen \\
\hline
8 & Motivation \& Relevanz & 10\% & 9/10 & Alltagsbezug (Kleidung beschreiben) \\
\hline
9 & Inklusion (UDL) & 10\% & 8/10 & Visuell + auditiv + schriftlich \\
\hline
10 & Praktikabilität & 5\% & 10/10 & 45 Min realistisch, Material vorhanden \\
\hline
\multicolumn{3}{|r|}{\textbf{GESAMT (gewichtet):}} & \textbf{9.1/10} & \\
\hline
\end{tabular}

\vspace{1em}
\textbf{Bewertung:}
\begin{itemize}
    \item[$\boxtimes$] \textcolor{successgreen}{$\geq$ 9.0: \textbf{FREIGABE} (Premium-Qualität)}
    \item[$\Box$] \textcolor{warningorange}{8.0--8.9: Iteration empfohlen}
    \item[$\Box$] 6.0--7.9: Überarbeitung erforderlich
    \item[$\Box$] $<$ 6.0: Grundlegende Neukonzeption
\end{itemize}

\vspace{1em}
\textbf{Stärken:}
\begin{itemize}
    \item Klare Struktur für 45-Minuten-Einzelstunde
    \item Differenzierung qualitativ (nicht nur quantitativ)
    \item Partnerübung erst nach Vokabelrevision (wie gewünscht)
    \item Praktikabel und zeitlich realistisch
\end{itemize}

\textbf{Mögliche Verbesserungen:}
\begin{itemize}
    \item Feedback-Qualität: Individuelle Rückmeldung bei EA verstärken
    \item UDL: Bildkarten könnten visuellen Zugang verstärken
\end{itemize}
\end{scorebox}

\end{document}
